\documentclass[../../../include/open-logic-section]{subfiles}

\begin{document}

\olfileid{sth}{z}{sep}
\olsection{وکھ کرن دا اصول}

سادہ خاصیت توں سیٹ بناؤن دے اصول دی تھاں اسیں ایس اصول توں شروع کردے آں:

\begin{axiom}[وکھ کرن دا خاکہ] ہر فارمولے $\phi(x)$ لئی ایہ اک مسلّمہ
اے کہ کسے وی $A$ لئی سیٹ $\Setabs{x \in A}{\phi(x)}$ موجود اے۔
\end{axiom}

دھیان دیو کہ ایہ اکّو مسلّمہ نہیں۔ ایہ مسلّمیاں دا اک \emph{خاکہ} اے۔
وکھ کرن دے \emph{لامتناہی} مسلّمے نیں، ہر فارمولے $\phi(x)$ لئی اک۔
خاکہ ایویں وی لکھ سکدے آں (تے عام طور اُتے ایویں ای لکھدے آں):

\begin{defish}
ہر اوہ فارمولا $\phi(x)$ جیہدے وچ ’’$S$‘‘ نہیں آؤندا، اوہدے لئی ایہ اک
مسلّمہ اے:
\[
	\forall A \exists S \forall x(x \in S \liff (\phi(x) \land x \in A)).
\]
\end{defish}

\olref[sth][][]{part} دے شروع وچ دسے رواج مطابق، وکھ کرن دے مسلّمیاں
وچ فارمولیاں~$\phi$ دے پیرامیٹر ہو سکدے نیں۔\footnote{ایہدا مطلب
سمجھن لئی \olref[sfr][infinite][induction]{natinductionschema} توں
جھٹ پچھوں والی گل بات ویکھو۔}

سیٹاں دا اکٹھا کرن والا مرحلہ وار تصور، جیہڑا اسیں بیان کردے آئے آں،
وکھ کرن دا جواز سدھا دیندا اے۔ ویکھن لئی منّ لو $A$ اک سیٹ اے۔
سو $A$ کسے مرحلے~$S$ اُتے بندا اے (\stageshier مطابق)۔ چونکہ $A$
مرحلے~$S$ اُتے بنیا سی، ایس لئی $A$ دے سارے رکن مرحلے $S$ توں
پہلاں بنے سن (\stagesacc مطابق)۔ ہُن خاص کر کے اوہ سارے سیٹ ویکھو
جیہڑے $A$ دے رکن وی نیں تے $\phi$ نوں پورا وی کردے نیں۔ صاف اے
کہ ایہ سارے سیٹ وی مرحلے~$S$ توں پہلاں بنے سن۔ ایس لئی ایہ وی اک سیٹ
$\Setabs{x \in A}{\phi(x)}$ بناندے نیں، مرحلے~$S$ اُتے (\stagesacc مطابق)۔

سادہ خاصیت توں سیٹ بناؤن دے اصول دے اُلٹ، ایہ رسل دے پیراڈوکس توں
بچدا اے۔ اسیں سیٹ $\Setabs{x}{x \notin x}$ دی موجودگی سدھی نہیں آکھ
سکدے۔ پر جے سانوں کوئی سیٹ~$A$ \emph{دتا ہووے}، تاں اسیں سیٹ
$R_A = \Setabs{x \in A}{x \notin x}$ دی موجودگی آکھ سکدے آں۔ پر ایس توں
اینا ای ثابت ہوندا اے کہ $R_A \notin R_A$ تے $R_A \notin A$۔ ایہناں وچ
کوئی بہتی اوکڑ والی گل نہیں۔

پر وکھ کرن دا اک سدھا تے چونکا دین والا نتیجہ اے:

\begin{thm}\ollabel{thm:NoUniversalSet}
کوئی \emph{ہمہ گیر} سیٹ نہیں، یعنی $\Setabs{x}{x = x}$ موجود نہیں۔
\end{thm}

\begin{proof}
تضاد کڈھن لئی منّ لو $V$ اک ہمہ گیر سیٹ اے۔ تاں وکھ کرن دے اصول نال
$R = \Setabs{x \in V}{x \notin x} = \Setabs{x}{x \notin x}$ موجود اے۔
ایہ رسل دے پیراڈوکس دے اُلٹ اے۔
\end{proof}

ہمہ گیر سیٹ دا نہ ہونا، بلکہ سیٹاں دی درجہ بندی دا کوئی آخری سرا نہ ہونا،
اکٹھا کرن والے مرحلہ وار تصور دیاں سبھ توں بنیادی گلاں وچوں اک اے۔
ایس لئی ایہ ویکھنا چنگا اے کہ سمجھ مطابق اسیں ہور راہ توں وی ایس گل
تک پہنچ سکدے آں۔ ہمہ گیر سیٹ لازماً اپنے آپ دا !!a{element} ہووے گا۔
پر ساڈے اکٹھا کرن والے مرحلہ وار تصور وچ ہر سیٹ درجہ بندی وچ پہلی واری
اوس پہلے مرحلے اُتے آؤندا اے جیہڑا اوہدے سارے !!{element}s توں جھٹ
پچھوں اے۔ ایس توں نکلدا اے کہ \emph{کوئی وی} سیٹ اپنے آپ دا رکن نہیں۔
اپنے آپ دا رکن سیٹ پہلی واری اوس مرحلے توں جھٹ پچھوں آؤنا سی جیہدے
اُتے اوہ پہلی واری آیا، تے ایہ بے تُکی گل اے۔ (\olref[sth][spine][rank]{defnsetrank}
وچ ویکھاں گے کہ سیٹ دے ’’درجے‘‘ دا خیال ورت کے ایس گل نوں ہور پکا
کیویں کر سکدے آں۔ پر ایہ کرن لئی سانوں پہلاں کجھ ہور مسلّمے چاہیدے ہون گے۔)

وکھ کرن تے عنصراں نال تعیّن دے کجھ ہور نتیجے ہِیٹھاں نیں۔

\begin{prop}\ollabel{prop:emptyexists}
جے کوئی سیٹ موجود اے، تاں $\emptyset$ موجود اے۔
\end{prop}

\begin{proof}
جے $A$ اک سیٹ اے، تاں $\emptyset = \Setabs{x \in A}{x \neq x}$
وکھ کرن دے اصول نال موجود اے۔
\end{proof}

\begin{prop}
$A \setminus B$ کسے وی دو سیٹاں $A$ تے $B$ لئی موجود اے۔
\end{prop}

\begin{proof}
$A \setminus B = \Setabs{x \in A}{x \notin B}$ وکھ کرن دے اصول نال موجود اے۔
\end{proof}

ایہ وی پتا لگدا اے کہ (تقریباً) کسے وی خاندان دے اشتراک موجود نیں:

\begin{prop}\ollabel{prop:intersectionsexist}
جے $A \neq \emptyset$، تاں $\bigcap A = \Setabs{x}{(\forall y \in A)x \in y}$
موجود اے۔
\end{prop}

\begin{proof}
منّ لو $A \neq \emptyset$، سو کوئی $c \in A$ اے۔ فیر $\bigcap A =
\Setabs{x}{(\forall y \in A)x \in y} = \Setabs{x \in c}{(\forall y \in
A)x \in y}$، جیہڑا وکھ کرن دے اصول نال موجود اے۔
\end{proof}

پر شرط $A \neq \emptyset$ دا دھیان دیو۔ کیوں جے $\bigcap \emptyset$
خالی شرط کارن ہمہ گیر سیٹ ہووے گا، جیہڑا \olref{thm:NoUniversalSet}
دے اُلٹ اے۔

\end{document}
